\chapter{Model Description}

In this project, we use the BTW Sandpile Model to study self-organized criticality (SOC). The model uses a finite two-dimensional lattice to perform a sandpile system, and each lattice site corresponds to a local sand pile. In this chapter, we will introduce a detailed mathematical definition for the sandpile model, including the local rules. 

\section{The BTW Sandpile Model}

The two-dimensional lattice is defined as:

$$
L=\{1,\ldots,N\}\times\{1,\ldots,M\}
$$

where:

\begin{itemize}
    \item $N$ is the number of rows;
    \item $M$ is the number of columns;
    \item $(i,j)\in L$ represents a lattice site.
\end{itemize}

At time step $t$, the number of sand grains at site $(i,j)$ is given by:

$$
p_t(i,j)\in\mathbb{N}_0. 
$$

The full sandpile configuration at time $t$ is written as:

$$
P_t=\{p_t(i,j)\}_{(i,j)\in L}. 
$$

In this model, $t$ represents the external sand addition step, rather than the number of internal topplings during an avalanche. At each time step, one grain of sand is added to the system, while the internal toppling process is considered as instantaneous \citep{Project2020}. 



\section{Toppling Dynamics and Avalanche Process}

At each time step $t$, we randomly select a lattice site $(i_t,j_t)\in L$, and then add one grain of sand on it:

$$
p_t(i_t,j_t)\leftarrow p_t(i_t,j_t)+1. 
$$

The model uses a fixed stability threshold:

$$
k=4. 
$$

A site becomes unstable when:

$$
p_t(i,j)\ge4. 
$$

When a site becomes unstable, a toppling occurs on it. During one toppling process, the site loses four grains of sand, and one grain is transferred to each of its four nearest neighbours \citep{Project2020}:

$$
p_t(i,j)\leftarrow p_t(i,j)-4, 
$$

$$
p_t(i\pm1,j)\leftarrow p_t(i\pm1,j)+1, 
$$

$$
p_t(i,j\pm1)\leftarrow p_t(i,j\pm1)+1. 
$$

Note that we only consider the nearest neighbours, without considering diagonal neighbours \citep{Project2020}. The nearest neighbour set is defined as:

$$
\mathcal{N}(i,j)
=
\{
(i+1,j),
(i-1,j),
(i,j+1),
(i,j-1)
\}
$$

In this project, we use the open boundary condition. If a toppling happens at the boundary of the lattice, some grains may be removed outside the lattice and system. 

A single grain addition may cause a chain reaction of topplings. When one unstable site topples, its neighbouring sites may also become unstable and then begin toppling. This process is called an avalanche. In this project, an avalanche is defined as the full sequence of topplings caused by one grain addition until the system becomes stable \citep{Project2020}, which is given by: 

$$
p_t(i,j)<k,
\quad
\forall(i,j)\in L. 
$$



\section{Avalanche Properties}

In order to study the avalanche behaviour, this project considers four avalanche properties: topples, area, loss, and radius \citep{Project2020}. 

We denote the avalanche topples as $s$, which is the total number of topplings during one avalanche. Note that if one site topples multiple times during the same avalanche, then every toppling is counted in $s$; We denote the avalanche area as $a$, which is the number of unique lattice sites that topple during one avalanche; We denote the avalanche loss as $l$, which is the total number of sand grains that leave the lattice during one avalanche. 

We denote the avalanche length as $r$, which is defined as the spatial radius:

$$
r
=
\max_{x\in A}
\|x-x_0\|_2
$$

where:

\begin{itemize}
    \item $x_0$ is the dropping site of the avalanche;
    \item $A$ is the set of all toppled lattice sites;
    \item $\|\cdot\|_2$ represents the Euclidean distance.
\end{itemize}

This length measures the maximum Euclidean distance between the dropping site and all toppled sites, which represents how far the avalanche spreads on grid.

After each avalanche, the simulation records the corresponding values of avalanche properties as $(s,a,l,r)$, which is used for the following statistical analysis. 



\section{Example of an Avalanche Process}

In order to clearly illustrate the toppling dynamics and demonstrate how to calculate avalanche properties in the BTW sandpile model, we use a $3\times3$ lattice avalanche example:  

$$
\begin{bmatrix}
1 & 3 & 3\\
2 & 4 & 2\\
2 & 1 & 1
\end{bmatrix}
\rightarrow
\begin{bmatrix}
1 & 4 & 3\\
3 & 0 & 3\\
2 & 2 & 1
\end{bmatrix}
\rightarrow
\begin{bmatrix}
2 & 0 & 4\\
3 & 1 & 3\\
2 & 2 & 1
\end{bmatrix}
\rightarrow
\begin{bmatrix}
2 & 1 & 0\\
3 & 1 & 4\\
2 & 2 & 1
\end{bmatrix}
\rightarrow
\begin{bmatrix}
2 & 1 & 1\\
3 & 2 & 0\\
2 & 2 & 2
\end{bmatrix}
$$

Here we assume the dropping site is $(2, 2)$. The avalanche process of this example is derived by implementing the definition and rules as shown above. Then we calculate the avalanche properties by following the definition. 

First, the avalanche topples $s$ denote the total number of toppling events during the avalanche. In this example, four topplings happen at the following lattice sites:

$$
(2,2), \quad (1,2), \quad (1,3), \quad (2,3), 
$$

then, 

$$
s=4. 
$$

Next, the avalanche area $a$ denotes the number of unique lattice sites that topple during the avalanche. Since all four toppled sites are different:

$$
a=4. 
$$

The avalanche loss $l$ denotes the number of sand grains that leave the system through the lattice boundary. In this example:

\begin{itemize}
    \item the central site $(2,2)$ is an interior site and produces no loss;
    \item the upper boundary site $(1,2)$ has one missing neighbour and produces one grain loss;
    \item the upper-right corner site $(1,3)$ has two missing neighbours and produces two grain losses;
    \item the right boundary site $(2,3)$ has one missing neighbour and produces one grain loss.
\end{itemize}

Then, the total loss is:

$$
l=0+1+2+1=4
$$

Finally, the avalanche radius $r$ is defined as the maximum Euclidean distance between the dropping site and all toppled sites:

$$
r
=
\max_{x\in A}
\|x-x_0\|_2
$$

where the dropping site is $x_0=(2,2)$, and the set of toppled sites is $A=\{(2,2),(1,2),(1,3),(2,3)\}$. 

Then the corresponding Euclidean distances are $0, 1, 1, \sqrt2$, and: 

$$
r=\sqrt2
$$

Hence, the avalanche observables for this example are:

$$
s=4,\quad a=4,\quad l=4,\quad r=\sqrt2. 
$$